\cotefile{53}
\scanpage{23}
% loc 1–2, opening bracket
[\unsure{this|cel} expresses that the homomorphism $A\mathrel{\tikz[baseline=-.3ex,x=1em,y=1em]{\draw[line width=.45pt] (0,0)--(1,0)--(1,.17);}}k'$ is a homomorphism of algebras]

% loc 3–6, Si G
If $G$ is endowed with a law of composition likewise $\mathrm{Dist}(D)$, the latter is associative (commutative) if and only if that of $G$ is, and \deleted{unital} $G$ admits a two-sided unit if and only if $\mathrm{Dist}(D)$ has a unit \deleted{\unsure{element}} $e$ such that $\delta(e)=e\boxtimes e$.

% loc 7–9, Dire que
To say that $G$ acts on a module $M$ \inserted{on the left)} means that $\mathrm{Dist}(G)$ acts on it on the left -- and this action extends that of $G$.

% loc 10–14, NB
\uline{NB} In certain cases, $\mathrm{Dist}(X)$ is reconstructed from $\mathrm{dist}(X,k)=\mathrm{Dst}(X)(k)$, for example if $A$ is \deleted{\unsure{finite type|of finite type}} projective of finite type over $k$, or more generally if $A$ is a sum of projective modules of finite type (for example \deleted{\unsure{if}} $k$ a field) -- but then one must consider $\mathrm{Dist}(x,k)$ as a topological module.

% loc 15–20, Re NB
\uline{Re NB} Especially interesting for monoid. In addition \deleted{\unsure{to|to the}} \inserted{to the subfunctor} $G$, $\mathrm{dist}(G)$ ' [resp $\mathrm{Dist}(G,k)$] contains all the infinitesimal information of $G$ : for example the \uline{Lie algebra} of $G$ is contained in $\mathrm{Dist}(G,k)$, -- \deleted{\unsure{and the operations}} its $[\,]$ law is induced by the bracket of the algebra, and in positive characteristic its p-th power too -- as well as the hyperalgebra.

\begin{center}\rule{1.6cm}{0.55pt}\end{center}
% loc 21, between the two rules
Case of constant schemes : one recovers the \uline{algebra} of a finite group.

\begin{center}\rule{1.6cm}{0.55pt}\end{center}
% loc 22, second Cas
Case of the diagonalizable group
% loc 23, A = k(M)
\[
 A=k(M)\quad\text{group algebra with coefficients in } k
\]
% loc 24–25, two paired rows
\diagram{\begin{tikzcd}[ampersand replacement=\&,column sep=small,row sep=large]
 A\otimes A \arrow[r,dash] \& A \& \text{comes from} \& M\times M \arrow[r] \& M \& \text{group law}\\
 A\otimes A \& \underset{\tikz[baseline,x=1em,y=1em]{\draw[gray,line width=.35pt] (.18,.12)--(.16,.02);\draw[gray,line width=.35pt] (.10,-.06) to[out=-90,in=0] (-.12,-.32);\draw[gray,line width=.35pt] (.16,-.22) to[out=-90,in=0] (-.15,-.49);}}{A} \arrow[l] \& \text{---} \& M\times M \& M \arrow[l,swap,"\substack{\text{diagonal}\\\text{map}}"] \&
\end{tikzcd}}
% loc 26–27, lowest two lines
\[
\begin{aligned}
 &\mathrm{Dist}(\mathbb D(M))(k)=\deleted{\unsure{*}}k^{M}\quad\text{with the \uline{product} \uline{algebra} law}\\
 &\qquad\text{[Infinite linear combinations of characters]} \qquad
 \text{(\unsure{does not depend on the additive structure of|does not depend on the additive structure of} $M$)}
\end{aligned}
\]
% archive sheet number: 22
